\documentclass[11pt]{article}
\usepackage[T1]{fontenc}
\usepackage[utf8]{inputenc}
\usepackage{lmodern,amsmath,amssymb,amsthm,mathtools,microtype,booktabs,enumitem,xcolor}
\usepackage[margin=1in]{geometry}
\usepackage[colorlinks=true,linkcolor=blue!45!black,citecolor=blue!45!black,urlcolor=blue!45!black]{hyperref}
\newtheorem{theorem}{Theorem}[section]
\newtheorem{proposition}[theorem]{Proposition}
\newtheorem{lemma}[theorem]{Lemma}
\newtheorem{corollary}[theorem]{Corollary}
\theoremstyle{definition}
\newtheorem{definition}[theorem]{Definition}
\newtheorem{remark}[theorem]{Remark}
\newcommand{\R}{\mathbb R}
\newcommand{\E}{\mathbb E}
\newcommand{\norm}[1]{\left\lVert#1\right\rVert}
\setlength{\emergencystretch}{3em}
\setlist{itemsep=2pt,topsep=5pt}
\hypersetup{pdftitle={Width-uniform Bayesian generalization above the mean-field scaling},pdfauthor={Dmitry Vaintrob; Kaarel Hanni; OpenAI Codex}}
\title{Width-uniform Bayesian generalization\\above the mean-field scaling}
\author{Dmitry Vaintrob \and Kaarel H\"anni \and OpenAI Codex\thanks{AI assistant developed by OpenAI; contributed to mathematical analysis and drafting.}}
\date{October 15, 2026}
\begin{document}
\maketitle
\begin{abstract}
We study noiseless Bayesian learning by fully connected neural networks with bounded Lipschitz activations, a linear readout, and a Gaussian weight prior with scale $\sigma=N^{-\gamma/2}$. Before observing data the prior is conditioned on a fixed absolute margin at every input of the Boolean cube. For every fixed depth and every fixed $\gamma>1$, we prove a sample bound polynomial in the input dimension and the bounded-weight network size of the target, uniformly over every larger learner width. The result holds at the specified nonzero prior variance, rather than only in a small-variance limit at fixed width. A cylinder construction gives a common target-mass lower bound for all $\gamma\geq1$. Above the critical scaling, Gaussian isoperimetry controls the Frobenius mass of singular values outside the vanishing operator-norm bulk. This yields a bound on the total probability of functions that lack a small description. Combining this bound with compression and a finite-class posterior argument proves width-uniform generalization without an exponential-in-dimension width assumption. A separate elementary argument covers the remaining smaller widths. We also prove polynomial equivalence between normalized negative log prior mass and bounded-weight network complexity. The full-depth learning theorem proved here is for $\gamma>1$; a companion paper treats the tanh mean-field endpoint through a different compression argument.
\end{abstract}

\section{Introduction and result}
The probability assigned to a function by a neural-network weight prior reflects both the cost of its smallest representation and the volume of its other representations. In a sufficiently strong Gaussian prior, the first quantity is easier to isolate. Earlier small-variance arguments take the variance to zero at each fixed width. Here we instead specify the variance as a power of the width and prove a bound valid at every width, for every fixed exponent above the mean-field value.

Our result is the following. Fix the depth $L$ and $\gamma>1$. If a Boolean target has a width-$C$ implementation with weights bounded by $C$, then the Bayesian learner defined below has sample complexity polynomial in $d,C,1/\varepsilon$, and $\log(1/\delta)$, independently of its width $N\geq C$. The polynomial degree may depend on $L$ and $\gamma$. The conclusion concerns posterior label error, with a corresponding bound for posterior majority prediction. It makes no claim about the computation needed to sample the posterior.

The proof develops the argument in Vaintrob and H\"anni's lecture notes~\cite{notes} and the finite-$\gamma$ draft~\cite{fable}. It has three quantitative ingredients. First, a small target can be cloned into any larger network while constraining only linearly many Gaussian coordinates. Second, at $\gamma>1$, the total probability of functions whose implementation needs large weighted norm is exponentially small on the $N^\gamma$ scale. Third, bounded weighted norm implies a width-independent finite description. The second step controls a single exceptional parameter event, avoiding a union bound over all Boolean functions. We include the compression proof in an appendix so that the learning theorem is self-contained apart from standard Gaussian concentration and isoperimetry.

\section{Model and prediction convention}\label{sec:model}
Fix $d\geq1$, $L\geq2$, and $\mathcal X=\{-1,1\}^d$. Let $\bar x=(1,x_1,\ldots,x_d)$, and let $\phi:\R\to[-1,1]$ be $1$-Lipschitz. The architecture has $L-1$ hidden layers of common width $N$ and a linear readout:
\[
 h^0(x)=\bar x,\qquad h^s(x)=\phi(W^s h^{s-1}(x))\quad(1\leq s<L),
 \qquad z_w(x)=W^Lh^{L-1}(x).
\]
There are no additional biases beyond the constant input coordinate. Tanh satisfies the activation assumptions. Non-polynomiality is not needed for the theorem when the target is assumed representable.

For a Boolean function $g$, set
\[
 U_{N,g}=\{w:g(x)z_w(x)>1/2\text{ for every }x\in\mathcal X\},
 \qquad \mathcal H_N=\{g:U_{N,g}\ne\varnothing\},
 \qquad U_N=\bigcup_{g\in\mathcal H_N}U_{N,g}.
\]
Define
\[
 \kappa_N(w)=\frac dN\|W^1\|_F^2+
 \sum_{s=2}^{L-1}\|W^s\|_F^2+N\|W^L\|_2^2,
 \qquad \kappa_N(g)=\inf_{w\in U_{N,g}}\kappa_N(w).
\]
Let $C_L(g)$ be the smallest integer $C\geq1$ for which a width-$C$ network in $U_{C,g}$ has every weight bounded in absolute value by $C$. We assume $C_L(f)<\infty$ for the target. This is bounded-weight network complexity at fixed depth; it does not silently identify arbitrary Boolean circuit size with fixed-depth network size.

Under the raw law $\mu_{N,\gamma}$ all weights are independent centered Gaussians with standard deviations
\[
 \sigma_1=N^{(1-\gamma)/2}/\sqrt d,\qquad
 \sigma_s=N^{-\gamma/2}\ (2\leq s<L),\qquad
 \sigma_L=N^{-(\gamma+1)/2}.
\]
The density is proportional to $\exp[-N^\gamma\kappa_N(w)/2]$.
The function prior is the raw law conditioned on $U_N$:
\[
 \pi_{N,\gamma}(g)=\frac{\mu_{N,\gamma}(U_{N,g})}{\mu_{N,\gamma}(U_N)}.
\]
Given $S_n=((x_i,f(x_i)))_{i=1}^n$ with independent $x_i\sim D$, the posterior $\Pi(\cdot\mid S_n)$ restricts $\pi_{N,\gamma}$ to functions satisfying all observed labels. For any distribution $D$ on $\mathcal X$, write
\[
 \operatorname{err}_D(g)=\Pr_{x\sim D}\{g(x)\ne f(x)\},
 \qquad R_D(\Pi)=\E_{g\sim\Pi}\operatorname{err}_D(g).
\]
The global margin condition applies before training, including at unobserved inputs. It is part of the learner, and differs from imposing margins only on the training sample. The theorem controls the posterior probability of an incorrect label, averaged over a fresh input.

% Standalone section: macros deliberately limited to standard AMS commands.
\section{A cylinder construction valid throughout \texorpdfstring{$\gamma\geq1$}{gamma >= 1}}
\label{sec:cloning}
The following construction bounds the target's prior cost at every width $N\geq C$,
without a width threshold diverging as $\gamma\downarrow1$.
We replicate the neurons of a small target network and constrain selected
Gaussian coordinates to remain near that implementation. This defines a
\emph{cylinder event}: the remaining, orthogonal coordinates retain their
Gaussian distribution. Controlling their operator norms then ensures that
the entire network preserves the target's margin.

\begin{proposition}[Uniform cloning mass bound]
Fix $L\geq2$ and a bounded $1$-Lipschitz activation. There is a function $Q_L(d,C)$ bounded by a polynomial such that for every $\gamma\geq1$, every Boolean target with
$C_L(f)\leq C$, and every $N\geq C$,
\[
 -\frac{2}{N^\gamma}\log\Pr\{f(x)z_w(x)>1/2\text{ for all }x\}
 \leq Q_L(d,C).
\]
The prior variances and absence of later-layer biases are exactly those in the
main model.
\end{proposition}
\begin{proof}
Choose a width-$C$ teacher, increase its readout by a factor of three, and write
its hidden matrices as $B^s$ and readout as $v$. Thus the teacher has margin
strictly greater than $3/2$, hidden entries bounded by $C$, and readout entries
bounded by $3C$. If the original witness has smaller width, first use its actual
width in the construction below; all estimates only improve.
Partition $\{1,\ldots,N\}$ into $C$ groups of sizes $k_a$ differing by at most
one. Then
\[
 1\leq k_a,\qquad \frac{\max k_a}{\min k_a}\leq2,
 \qquad \frac N{k_a}\leq2C.
\]
Let $u_a=k_a^{-1/2}{\bf1}_{\text{group }a}$ and let $P$ be the orthogonal
projection onto their span. The exact clone $w_*$ has first-layer row $B^1_b$
for a neuron in group $b$, hidden entry $B^s_{ba}/k_a$ from group $a$ to group
$b$, and readout $v_b/k_b$ on group $b$. Every hidden activation of the exact
clone is group-constant. Moreover,
\[
 \|W_*^s\|_{\mathrm{op}}\leq\sqrt2 C^2\quad(2\leq s<L),
 \qquad \sqrt N\|W_*^L\|_2\leq3\sqrt2 C^2.
\]
The first bound follows by restricting to the group-constant subspace: the
resulting matrix has entries $\sqrt{k_b/k_a}B^s_{ba}$, and its Frobenius norm
is at most $\sqrt2 C^2$.

Set $t=N^{(1-\gamma)/2}\leq1$. The random coordinates
$\sqrt d W^1_{ij}$, $\sqrt N W^s_{ij}$ for $2\leq s<L$, and
$\alpha_j=NW^L_j$ are independent centered Gaussians of standard deviation
$t$. Impose the following cylinder event $\mathcal C_\rho$, for
$0<\rho\leq1$:
\begin{itemize}
\item every $\sqrt d W^1_{ij}$ is within $\rho t$ of
$\sqrt d (W_*^1)_{ij}$;
\item for every hidden row $j$ and group $a$, the coordinate
$\langle\sqrt N W^s_{j,\cdot},u_a\rangle$ is within $\rho t$ of its
cloned value $\sqrt{N/k_a}B^s_{ba}$;
\item every $\alpha_j$ is within $\rho t$ of its cloned value $Nv_b/k_b$.
\end{itemize}
All these constraints concern independent scalar Gaussian coordinates.
The orthogonal coordinates of each internal row remain independent prior
Gaussians. Consequently write
\[
 W^s=W_*^s+R^s+E^s,\qquad R^s=R^sP,\quad E^sP=0.
\]
On the cylinder, $\|R^s\|_{\mathrm{op}}\leq\|R^s\|_F\leq\rho t\sqrt C$.
The free matrix is distributed as $(t/\sqrt N)G^s(I-P)$, with $G^s$ a standard
Gaussian matrix, independently of the cylinder coordinates. Let
\[
 b_L=2+\sqrt{2\log(2L)}.
\]
The Gaussian operator-norm bound and a union bound imply, conditionally on
$\mathcal C_\rho$, that all $\|E^s\|_{\mathrm{op}}\leq b_Lt$ with probability
at least $1/2$. This is also valid when $L=2$, where there are no such matrices.

Let $e_s(x)=N^{-1/2}\|h^s_w(x)-h^s_*(x)\|_2$. The input constraints give
$e_1(x)\leq\rho(d+1)/\sqrt d\leq\rho(d+1)$. Since $E^sh^{s-1}_*=0$,
boundedness and Lipschitz continuity give, simultaneously for every input,
\[
 e_s(x)\leq (\sqrt2 C^2+\sqrt C+b_L)e_{s-1}(x)+\rho\sqrt C.
\]
Put $A=2+3C^2+b_L$. Induction yields
$e_s(x)\leq\rho(d+L+1)A^{L-2}$ for $s<L$. At the output,
\[
 |z_w(x)-z_*(x)|
 \leq3\sqrt2 C^2 e_{L-1}(x)+\rho.
\]
Hence choosing
\[
 \rho=\bigl[2+9C^2(d+L+1)A^{L-2}\bigr]^{-1}
\]
ensures output discrepancy at most $1/2$. All networks in the cylinder with
the operator-norm event therefore implement $f$ with margin greater than one.

It remains to price the cylinder. Its number of constrained coordinates is
\[
 m=N\bigl[d+2+(L-2)C\bigr].
\]
The sum of their squared centers is at most $NB_0$, where one may take
\[
 B_0=d(d+1)C^2+2(L-2)C^4+36C^4.
\]
Indeed, a first-layer center has magnitude at most $\sqrt d C$, an internal
center at most $\sqrt{2C}C$, and a readout center at most $6C^2$.
For $Z\sim N(0,t^2)$, every real $a$, and every $\rho>0$,
\[
 \Pr\{|Z-a|\leq\rho t\}
 \geq\frac{2\rho}{\sqrt{2\pi}}
       \exp\{-a^2/t^2-\rho^2\}.
\]
Multiplication and $t^{-2}=N^{\gamma-1}$ show
\[
 -\log\Pr(\mathcal C_\rho)
 \leq N^\gamma B_0
 +N[d+2+(L-2)C]\log\frac{e\sqrt{2\pi}}{2\rho}.
\]
The desired probability is at least $\Pr(\mathcal C_\rho)/2$. Since
$N\leq N^\gamma$ and $N^\gamma\geq1$, this proves the claim with
\[
 Q_L(d,C)=2B_0+\frac{4[d+2+(L-2)C]}{\rho}+2.
\]
Here we used $\log(e\sqrt{2\pi}/(2\rho))\leq2/\rho$ for $0<\rho\leq1$.
The displayed expression is a polynomial in $d,C$ at every fixed depth.
\end{proof}

% Paper-ready fragment. Depends only on the model definitions, the cylinder
% mass proposition in unified-cloning.tex, and the earlier compression theorem.
% Use \mathcal A_L for that proposition's mass bound (called Q_L there).

\section{Width-independent learning at every fixed $\gamma>1$}
Let $\mathcal A_L(d,C)=\max\{1,Q_L(d,C)\}$, where $Q_L$ is the bound
from the preceding cylinder construction. Then $\mathcal A_L$ is
polynomially bounded and
\begin{equation}\label{eq:uniform-target-mass}
 \mu_{N,\gamma}(U_{N,f})\geq
 \exp\{-\mathcal A_L(d,C)N^\gamma/2\}
 \qquad(\gamma\geq1,\ N\geq C).
\end{equation}
We now control the \emph{total} mass of functions of large circuit complexity.
A Gaussian geometric argument gives this control for every fixed $\gamma>1$.

\begin{lemma}[A Gaussian singular-value tail]\label{lem:gaussian-spectral-tail}
Let $M$ be an $N\times N$ matrix with independent $N(0,\sigma^2)$ entries.
For $a\geq8\sigma\sqrt N$, let $M_{>a}$ retain exactly the singular values of
$M$ greater than $a$. Then, for every $K>0$,
\[
 \Pr\{\|M_{>a}\|_F^2>K\}\leq
 \exp\{-K/(8\sigma^2)\}.
\]
\end{lemma}
\begin{proof}
The standard Gaussian operator-norm inequality~\cite{davidson} implies that
the ball $B_a=\{T:\|T\|_{\mathrm{op}}\leq a/2\}$ has Gaussian
probability at least $1-e^{-2N}>1/2$. Indeed,
\[
 \Pr\{\|M\|_{\mathrm{op}}>2\sigma\sqrt N+\sigma t\}
 \leq e^{-t^2/2},
\]
and $a/2\geq4\sigma\sqrt N$. The Frobenius distance to $B_a$ is
\[
 D(M)=\operatorname{dist}_F(M,B_a)
 =\left[\sum_j(s_j(M)-a/2)_+^2\right]^{1/2}.
\]
The Gaussian isoperimetric inequality~\cite{borell,ledoux}, applied to the Euclidean enlargement
of $B_a$ in the $N^2$ matrix coordinates, gives
\[
 \Pr\{D(M)>t\}\leq1-\Phi(t/\sigma)
 \leq e^{-t^2/(2\sigma^2)},\qquad t>0,
\]
where $\Phi$ is the standard normal distribution function. If $s>a$, then
$s\leq2(s-a/2)$, so $\|M_{>a}\|_F\leq2D(M)$. Apply the last bound with
$t=\sqrt K/2$.
\end{proof}

\begin{lemma}[Most margin-separated draws have low-norm descriptions]
\label{lem:aggregate-structural-tail}
Fix $K\geq1$, and set
\[
 a_K=\frac{1}{8L(\sqrt K+1)^{L-2}},\qquad
 J_K=(L+3)K.
\]
Suppose $\gamma>1$ and
\begin{equation}\label{eq:spectral-width-conditions}
 N^{\gamma-1}\geq8(d+1),\qquad
 N^{\gamma-1}\geq64a_K^{-2}.
\end{equation}
For
$\mathcal S_N(K)=\{g\in\mathcal H_N:\kappa_N(g)\leq J_K\}$,
\begin{equation}\label{eq:aggregate-tail}
 \mu_{N,\gamma}\left(\bigcup_{g\notin\mathcal S_N(K)}U_{N,g}\right)
 \leq L e^{-KN^\gamma/8}.
\end{equation}
\end{lemma}
\begin{proof}
Let $G_K$ be the event
\[
 \frac dN\|W^1\|_F^2\leq K,\qquad
 N\|W^L\|_2^2\leq K,\qquad
 \|(W^s)_{>a_K}\|_F^2\leq K\quad(2\leq s<L).
\]
The first and last weighted norms have distributions
$N^{-\gamma}\chi^2_{N(d+1)}$ and $N^{-\gamma}\chi^2_N$.
For $X\sim\chi_D^2$ and $u\geq8D$, the exponential-moment bound at
parameter $3/8$ gives
\[
 \Pr\{X>u\}\leq 4^{D/2}e^{-3u/8}\leq e^{-u/4}.
\]
The first condition in~\eqref{eq:spectral-width-conditions} therefore bounds
each of these two failure probabilities by $e^{-KN^\gamma/4}$.
For each internal layer, the second condition gives
$a_K\geq8N^{(1-\gamma)/2}$, so Lemma~\ref{lem:gaussian-spectral-tail}
bounds its failure probability by $e^{-KN^\gamma/8}$.
A union bound yields
\begin{equation}\label{eq:good-draw-tail}
 \mu_{N,\gamma}(G_K^c)\leq L e^{-KN^\gamma/8}.
\end{equation}

Fix a draw $w\in G_K$. Replace each internal matrix by its
$>a_K$ singular-value truncation, obtaining $\widetilde w$.
The retained matrices have operator norms at most $\sqrt K$; the discarded
matrices have operator norms at most $a_K$.
For any input, let $e_s$ be the Euclidean difference of the original and
truncated activations at hidden layer $s$. Boundedness and Lipschitz
continuity of the activation give
\[
 e_1=0,\qquad e_s\leq\sqrt K\,e_{s-1}+a_K\sqrt N
 \quad(2\leq s<L).
\]
The readout has norm at most $\sqrt{K/N}$. Thus, when $L\geq3$,
\[
 |z_{\widetilde w}(x)-z_w(x)|
 \leq a_K\sum_{j=1}^{L-2}K^{j/2}
 \leq a_K(L-2)K^{(L-2)/2}\leq\frac18.
\]
When $L=2$, the networks are identical, so the same final bound holds.
All these bounds hold simultaneously for every input.
If $w\in U_{N,g}$, the truncated network has strict margin $3/8$ for $g$.
Doubling its readout restores strict margin $1/2$ and gives weighted norm
at most $K+(L-2)K+4K=J_K$. Therefore every $w\in G_K\cap U_N$ implements
a function in $\mathcal S_N(K)$. The event on the left of
\eqref{eq:aggregate-tail} is contained in $G_K^c$, and
\eqref{eq:good-draw-tail} proves the claim.
\end{proof}

The key point is that~\eqref{eq:aggregate-tail} bounds an entire union of
function events directly. No union bound over the $2^{2^d}$ Boolean functions
is necessary. Conditioning the prior on $U_N$ does not affect its ratio to
the target event's mass.

\paragraph{The finite capacity bound used below.}
The compression theorem from the strong-prior argument gives a polynomially
bounded $m_L(d,J+1)$ such that every function with $\kappa_N(g)\leq J$ has
a width-$m=m_L(d,J+1)$ realization, with strict margin $1/2$ and all weights
bounded by $m$. The $+1$ permits an arbitrarily near minimizer when the
infimum defining $\kappa_N(g)$ is not attained. For definiteness one can take
\[
 \rho=\frac{2^{-L}(J+1)^{-(L-1)/2}}{12(L-1)},\qquad
 m_L(d,J+1)=\max\left\{\lceil9(J+1)\rceil,
 \left\lceil256\rho^{-4}\bigl(d+\log(L/\rho)\bigr)\right\rceil\right\}.
\]
Writing $P_m=m(d+2)+(L-2)m^2$, define
\begin{equation}\label{eq:finite-capacity-bound}
 H_L(d,J)=P_m\log\left(
 3+16Lm\max\{m,d+1\}[m(m+1)]^{L-1}\right).
\end{equation}
Then
\begin{equation}\label{eq:entropy-norm-class}
 \log\#\{g\in\mathcal H_N:\kappa_N(g)\leq J\}
 \leq H_L(d,J).
\end{equation}
To see this, round each weight of the compressed realization to a grid of
spacing
$r=[8L\max\{m,d+1\}\{m(m+1)\}^{L-1}]^{-1}$.
The rounding changes each weight by at most $r$ and keeps weights bounded
by $m+1$. A layerwise telescoping estimate in the maximum norm bounds the
output change by
$L\max\{m,d+1\}r[m(m+1)]^{L-1}=1/8$.
Thus every rounded array has the same Boolean signs as its original
realization. Distinct functions require distinct arrays, and the number of
arrays is at most $(3+2m/r)^{P_m}$. This proves
\eqref{eq:entropy-norm-class}. At fixed $L$, $H_L$ is bounded by a
polynomial in its arguments.

\begin{theorem}[All widths at fixed $\gamma>1$]
\label{thm:all-width-gamma-gt-one}
Fix $L\geq2$ and $\gamma>1$. There exists a polynomial $p_{L,\gamma}$ such
that the following holds. Let $C=C_L(f)<\infty$, let $N\geq C$, and let the
function prior and consistency posterior use the global margin conditioning
specified above. For every distribution $D$ on $\{-1,1\}^d$, every
$\varepsilon,\eta\in(0,1)$, and
\[
 n\geq\frac{2}{\varepsilon}\left[
 p_{L,\gamma}\bigl(d,C,\log(1/\varepsilon)\bigr)
 +\log\frac{2}{\varepsilon\eta}\right],
\]
with probability at least $1-\eta$ over $n$ independent labeled examples,
\[
 \mathbb E_{g\sim\Pi(\cdot\mid S_n)}\operatorname{err}_D(g)
 \leq\varepsilon.
\]
The polynomial is independent of $N$. Its exponent may depend on $\gamma$;
this theorem makes no uniform assertion as $\gamma\downarrow1$.
\end{theorem}
\begin{proof}
Set
\[
 A=\mathcal A_L(d,C),\qquad K=4A+8\log(2L/\varepsilon),
 \qquad a=a_K,
\]
and define the integer threshold
\[
 T=\left\lceil\max\left\{C,
 [8(d+1)]^{1/(\gamma-1)},
 (8/a)^{2/(\gamma-1)}\right\}\right\rceil.
\]
First suppose $N\geq T$. By the aggregate structural tail and the target
mass bound,
\begin{align*}
 \frac{\pi_{N,\gamma}(\mathcal H_N\setminus\mathcal S_N(K))}
 {\pi_{N,\gamma}(f)}
 &=\frac{\mu_{N,\gamma}(\bigcup_{g\notin\mathcal S_N(K)}U_{N,g})}
 {\mu_{N,\gamma}(U_{N,f})}\\
 &\leq L\exp\{-(K/8-A/2)N^\gamma\}\leq\varepsilon/2.
\end{align*}
The set $\mathcal S_N(K)$ contains $f$: otherwise the left side would be
at least one. Its log cardinality is at most
$H_L(d,J_K)$ by~\eqref{eq:entropy-norm-class}.
Any particular function of true error greater than $\varepsilon/2$ is
consistent with $n$ independent examples with probability at most
$e^{-n\varepsilon/2}$. A union bound shows that, if
\[
 n\geq\frac2\varepsilon
 \left[H_L(d,J_K)+\log(1/\eta)\right],
\]
then, with probability at least $1-\eta$, every consistent function in
$\mathcal S_N(K)$ has error at most $\varepsilon/2$.
On every training sample, the posterior denominator is at least
$\pi_{N,\gamma}(f)$, so posterior mass outside this set is at most
$\varepsilon/2$. The posterior risk is therefore at most $\varepsilon$.

Now suppose $C\leq N<T$. Since $\mu_{N,\gamma}(U_N)\leq1$,
\[
 -\log\pi_{N,\gamma}(f)
 \leq-\log\mu_{N,\gamma}(U_{N,f})
 \leq AN^\gamma/2\leq AT^\gamma/2.
\]
For any prior with positive target mass, let $B_n$ be posterior mass assigned
to functions with true error greater than $\varepsilon/2$. The posterior
denominator is at least the target's prior mass, giving
\[
 \mathbb E_{S_n}B_n
 \leq\frac{1}{\pi_{N,\gamma}(f)}
 \sum_{g:\operatorname{err}_D(g)>\varepsilon/2}
 \pi_{N,\gamma}(g)(1-\operatorname{err}_D(g))^n
 \leq\frac{e^{-n\varepsilon/2}}{\pi_{N,\gamma}(f)}.
\]
Markov's inequality gives $B_n\leq\varepsilon/2$ with probability at least
$1-\eta$ as soon as
\[
 n\geq\frac2\varepsilon\left[AT^\gamma/2+
 \log\frac{2}{\varepsilon\eta}\right].
\]
The remaining posterior mass contributes at most $\varepsilon/2$ to risk.

It is consequently enough to majorize
\[
 \max\{H_L(d,(L+3)K),\ AT^\gamma/2\}
\]
by $p_{L,\gamma}(d,C,\log(1/\varepsilon))$.
The cylinder bound makes $A$ polynomially bounded in $d,C$ at fixed $L$;
then $K$ is polynomially bounded in $d,C,\log(1/\varepsilon)$.
For fixed $L,\gamma>1$, so are $a^{-1}$, $T$, $T^\gamma$, and the last
display. This proves the stated uniform-in-width polynomial guarantee.
\end{proof}

For a confidence threshold $q\in(0,1)$, define
\[
 F_q(\Pi;f,D)=\Pr_{x\sim D}\{\Pi(g(x)=f(x)\mid S_n)\leq q\}.
\]
This counts an input as successful only when the posterior probability of its
correct label is strictly greater than $q$. Since the posterior probability
of an error is at least $1-q$ on the failure set,
\[
 F_q(\Pi;f,D)\leq\frac{R_D(\Pi)}{1-q}.
\]
In particular $F_{2/3}\leq3R_D(\Pi)$. Applying the theorem at accuracy
$\varepsilon/3$ proves the same width-independent polynomial sample bound
for having correct-label confidence greater than $2/3$ on at least a
$1-\varepsilon$ fraction of test inputs.

A posterior-majority classifier has risk at most twice the posterior average
risk, since every majority mistake places posterior mass at least $1/2$ on
the wrong label. Applying the theorem with accuracy $\varepsilon/2$ therefore
gives the same type of guarantee for majority prediction.

\paragraph{Comparison with the supplied draft.}
The supplied draft's structural estimates are compatible with the aggregate
argument above, but its learning section instead sums a per-function bound
over all Boolean functions. That step introduces an unnecessary
$N^\gamma\gtrsim2^d$ condition. Controlling the bad parameter event directly
removes this exponential width requirement. The present theorem also covers
the remaining polynomial range of widths through the elementary posterior
bound. Neither this repair nor the spectral lemma establishes the endpoint
$\gamma=1$: at that endpoint the Gaussian bulk operator norm remains of
constant order, whereas the truncation error requires a small cutoff.


\section{Prior mass and bounded-weight network complexity}
The same estimates also give the mass-complexity equivalence proposed in
the supplied draft, now without an auxiliary width proviso.

\begin{corollary}[Polynomial equivalence of mass and network complexity]
\label{cor:mass-complexity}
For every fixed $L\geq2$ and $\gamma>1$, there is a polynomial
$P_{L,\gamma}$ with the following property. Define the unconditioned mass
complexity
\[
 \overline\kappa_{N,\gamma}(f)
 =-\frac{2}{N^\gamma}\log\mu_{N,\gamma}(U_{N,f}).
\]
At every $N\geq1$, whenever this quantity is finite,
\[
 C_L(f)\leq
 P_{L,\gamma}\bigl(d,1+\overline\kappa_{N,\gamma}(f)\bigr).
\]
Conversely, whenever $N\geq C=C_L(f)$,
\[
 \overline\kappa_{N,\gamma}(f)\leq Q_L(d,C).
\]
The converse holds for every $\gamma\geq1$, with the same polynomial
$Q_L$ from the cylinder construction. The first polynomial is independent
of $N$, but may have degree depending on $\gamma$.
\end{corollary}
\begin{proof}
The converse is the cylinder mass bound. For the first direction, write
$b=\overline\kappa_{N,\gamma}(f)$, $B=\max\{1,b\}$, and
$K=4B+8\log(2L)$. Set $a=[8L(\sqrt K+1)^{L-2}]^{-1}$ and
\[
 T=\left\lceil\max\{[8(d+1)]^{1/(\gamma-1)},
                         (8/a)^{2/(\gamma-1)}\}\right\rceil.
\]
This is polynomially bounded in $d,B$ at fixed $L,\gamma$.
If $N\geq T$, Lemma~\ref{lem:aggregate-structural-tail} gives
\[
 \mu_{N,\gamma}(G_K^c)
 \leq Le^{-KN^\gamma/8}
 \leq\tfrac12e^{-bN^\gamma/2}
 <\mu_{N,\gamma}(U_{N,f}).
\]
Thus some implementing draw is in $G_K$, and its truncated, rescaled
network implements $f$ with weighted norm at most $(L+3)K$.
Compression gives $C_L(f)\leq m_L(d,(L+3)K+1)$.

If $N<T$, let $D_N=N(d+2)+(L-2)N^2$.
Under the raw prior, $q=N^\gamma\kappa_N(w)\sim\chi^2_{D_N}$.
For $u=4bN^\gamma+2D_N\log2+4$, the exponential-moment bound gives
\[
 \Pr\{q\geq u\}\leq2^{D_N/2}e^{-u/4}
 =e^{-bN^\gamma-1}<\mu_{N,\gamma}(U_{N,f}).
\]
Consequently an implementing draw has
\[
 \kappa_N(w)<u/N^\gamma
 \leq4B+2\bigl[T(d+2)+(L-2)T^2\bigr]\log2+4=:H.
\]
Another application of compression gives $C_L(f)\leq m_L(d,H+1)$.
Both resulting bounds are polynomially bounded in $d,B$ at fixed
$L,\gamma$, proving the claim.
\end{proof}


\section{Scope of the result}
The theorem supplies an explicit finite-variance extension of the strong-prior argument. For large enough width the relevant-class sample bound depends on $L,d,C,\log(1/\varepsilon)$ but not on $\gamma$; $\gamma$ determines when the spectral cutoff lies above the Gaussian bulk. Covering all widths uses a polynomial whose degree grows as $\gamma$ approaches one. This is a limitation of the present argument, rather than a lower bound asserting that such deterioration is necessary.

The cloning bound is uniform over $\gamma\geq1$, including the endpoint. At $\gamma=1$ the internal Gaussian matrices have a bulk of constant operator norm, so the small-cutoff argument does not establish a full-depth compression theorem. The companion mean-field paper~\cite{critical} gives a different proof for tanh at arbitrary fixed depth: it compresses the quadratic-tilt representations of exact rate minimizers and obtains polynomial sample complexity for all sufficiently large widths. Its width threshold is not quantified. For one hidden layer, a direct chi-square argument gives the stronger every-width conclusion for all $\gamma\geq1$ and any bounded $1$-Lipschitz activation.

For $\gamma<1$, tanh's diverging preactivation variance leads to a saturated, iterated-arcsine kernel. Under global margin conditioning its learning limit depends further on the sign of $\gamma$. The companion kernel note~\cite{subcritical} treats these limits and explains why a lower bound for regression using the posterior mean cannot be substituted for a classification lower bound under the present conventions.

\appendix
\section{A complete norm-compression proof}\label{sec:compression}
The next theorem is deterministic. It applies at any width and does not refer to the prior exponent. It is the compression result used in the finite capacity bound above.
\begin{theorem}[Compression]\label{thm:compression}
Let $d\geq1$, $L\geq2$, and let $\phi:\mathbb R\to[-1,1]$ be
$1$-Lipschitz. Suppose a width-$N$ network $W$ satisfies
$\kappa_N(W)\leq K$, where $K\geq1$. Set
\begin{equation}\label{eq:compression-parameters}
 \begin{gathered}
 \eta=\frac{1}{12(L-1)},\qquad
 \rho=\eta 2^{-L}K^{-(L-1)/2},\\
 m_L(d,K):=M=\max\left\{\lceil9K\rceil,
 \left\lceil256\rho^{-4}\bigl(d+\log(L/\rho)\bigr)\right\rceil\right\}.
 \end{gathered}
\end{equation}
There is a width-$M$ network $\widetilde W$ such that
\begin{equation}\label{eq:compression-conclusion}
 \max_{x\in\{-1,1\}^d}
 |z_{\widetilde W}(x)-z_W(x)|\leq\frac1{12},
 \qquad \kappa_M(\widetilde W)\leq4K.
\end{equation}
If $W$ implements a Boolean function $g$ with strict margin $1/2$,
then multiplying the readout of $\widetilde W$ by $3/2$ gives a
width-$M$ implementation of $g$ with strict margin $1/2$,
weighted norm at most $9K$, and every weight of absolute value at most
$M$. In particular, $C_L(g)\leq M$. For fixed $L$, $M$ is bounded by a
polynomial in $d$ and $K$.
\end{theorem}

\begin{proof}
We first truncate large entries and then sample the hidden neurons.
All norms of rows below are Euclidean, and all matrix norms with
subscript $F$ are Frobenius norms.

\emph{Truncation.}
For $2\leq s\leq L$, let $r_i^s=\|W^s_{i,*}\|_2$. In each nonzero
row, delete entries larger in absolute value than
\[
 \tau_i^s=\frac{2r_i^s}{\rho\sqrt N}.
\]
Leave zero rows unchanged, write $\overline W^s$ for the resulting
matrices, and set $\overline W^1=W^1$. Here and below $W^L$ is the
single readout row. Because all hidden activations have absolute value
at most one, deleting a row's large entries changes its preactivation
by at most
\[
 \sum_{j:\,|W^s_{ij}|>\tau_i^s}|W^s_{ij}|
 \leq \frac{(r_i^s)^2}{\tau_i^s}
 =\frac\rho2\sqrt N\,r_i^s.
\]
The vector of changes at an internal layer therefore has norm at most
$(\rho/2)\sqrt N\|W^s\|_F$. Every original or truncated internal
matrix has operator norm at most $\sqrt K$, and the readout has norm
at most $\sqrt{K/N}$. The activation is a contraction in Euclidean
norm. Thus truncating layer $s<L$ changes the output by at most
\[
 \frac\rho2\sqrt{NK}\,K^{(L-s-1)/2}\sqrt{K/N}
 =\frac\rho2 K^{(L-s+1)/2}\leq\frac\eta2.
\]
Truncating the readout changes it by at most
$(\rho/2)\sqrt K\leq\eta/2$. These bounds also hold when other
layers have already been truncated. Telescoping over the $L-1$
truncations gives total output error at most $1/24$.

\emph{Independent samples and three events.}
For each hidden layer $s$, independently choose a list
$I_s(1),\ldots,I_s(M)$ of independent uniform indices in
$\{1,\ldots,N\}$. Repeated indices represent separate copies of a
neuron. Define the fully sampled network by
\begin{align*}
 \widetilde W^1_{a,*}&=W^1_{I_1(a),*},\\
 \widetilde W^s_{a,b}&=\frac NM
      \overline W^s_{I_s(a),I_{s-1}(b)}
      &&(2\leq s\leq L-1),\\
 \widetilde W^L_b&=\frac NM\overline W^L_{I_{L-1}(b)}.
\end{align*}
For the analysis expose these independent lists in reverse order,
$I_{L-1},\ldots,I_1$. At the stage sampling layer $m$, all earlier
layers still have their original, truncated weights. Let
$\overline h_j^m(x)$ denote their exactly computed activations. The
estimated sum into a retained row $i$ of layer $m+1$ is
\[
 \frac NM\sum_{b=1}^M
 \overline W^{m+1}_{i,I_m(b)}\overline h^m_{I_m(b)}(x).
\]
Its expectation is the corresponding full sum over $N$ neurons.
Conditional on the later lists, its summands are independent and
bounded in absolute value by $(N/M)\tau_i^{m+1}$. Hoeffding's inequality~\cite{hoeffding}
shows that its deviation exceeds
$(\rho/2)\sqrt N r_i^{m+1}$ with probability at most
\[
 2\exp\left(-\frac{\rho^4 M}{32}\right).
\]
A zero row is estimated exactly. The bound is uniform in the retained
row, so it remains valid after averaging over the later lists.

Let $A$ be the event that all these row estimates, at all stages and
all $2^d$ inputs, satisfy their deviation bounds. At the last hidden
layer there is one output row; at every other stage there are $M$
retained rows. Consequently
\begin{equation}\label{eq:compression-hoeffding}
 \Pr(A^c)\leq 2^{d+1}LM\exp(-\rho^4M/32)\leq\frac1{100}.
\end{equation}
For completeness, put $a=d+\log L+\log(1/\rho)$ and
$M_0=\lceil256\rho^{-4}a\rceil$. Then
$M_0\leq257\rho^{-4}a$, and $u\exp(-\rho^4u/32)$ decreases for
$u\geq M_0$. Substituting $M_0$ and using $\log a\leq a-1$ bounds
the logarithm of the right side before its final inequality by
\[
 -(7-\log2)d-6\log L-3\log(1/\rho)+\log514-1
 \leq-\log100.
\]
This also proves the bound when $M$ was enlarged to be at least $9K$.

Next define the total retained row mass using the \emph{original}
row norms:
\[
 R=N\sum_{s=2}^{L-1}\sum_{a=1}^M(r^s_{I_s(a)})^2
      +N\|W^L\|_2^2.
\]
Its expectation is
\[
 \mathbb E R
 =M\sum_{s=2}^{L-1}\|W^s\|_F^2+N\|W^L\|_2^2
 \leq MK.
\]
By Markov's inequality the event $B=\{R\leq8MK\}$ has probability
at least $7/8$. In particular, on $B$ the retained row mass at each
individual stage is at most $8MK$.

Finally, direct calculation gives
\begin{align*}
 \mathbb E\left[\frac dM\|\widetilde W^1\|_F^2\right]
     &=\frac dN\|W^1\|_F^2,\\
 \mathbb E\|\widetilde W^s\|_F^2
     &=\|\overline W^s\|_F^2 &&(2\leq s\leq L-1),\\
 \mathbb E\left[M\|\widetilde W^L\|_2^2\right]
     &=N\|\overline W^L\|_2^2.
\end{align*}
For an internal matrix, independence of the two adjacent layers'
index lists makes each index pair uniform on $\{1,\ldots,N\}^2$;
the factor $(N/M)^2$ is canceled by the $M^2/N^2$ sampling factor.
Independence among different matrix entries is unnecessary.
Thus $\mathbb E\kappa_M(\widetilde W)\leq K$, and the event
$C=\{\kappa_M(\widetilde W)\leq4K\}$ has probability at least $3/4$.
On $C$ every sampled internal matrix has operator norm at most
$2\sqrt K$, and the sampled readout has norm at most $2\sqrt{K/M}$.

The events need not be independent:
\[
 \Pr(A\cap B\cap C)\geq1-\frac1{100}-\frac18-\frac14>0.
\]
Fix a collection of lists in their intersection.

\emph{Reverse-order telescoping.}
Let $F_m$ be the network with hidden layers $m,\ldots,L-1$ sampled
and layers $1,\ldots,m-1$ still original and truncated. Incoming
weights to a sampled layer select its rows, and outgoing weights from
a sampled layer are multiplied by $N/M$. This defines the intermediate
networks with different widths in different layers. In particular,
$F_L$ is the truncated network and $F_1$ is the fully sampled network.

When comparing $F_{m+1}$ and $F_m$, layer $m$ computes the original
truncated neuron functions before selection. The only new discrepancy
is in the sums into the retained neurons of layer $m+1$.
For $m=L-1$, event $A$ bounds the output discrepancy directly by
$(\rho/2)\sqrt K\leq\eta/2$. For $m<L-1$, put $t=L-m$.
Events $A$ and $B$ bound the Euclidean discrepancy in layer $m+1$'s
preactivation by $(\rho/2)\sqrt{8MK}$. All matrices \emph{after}
layer $m+1$ already equal their final sampled matrices, so their
operator norms are controlled by $C$. Propagating the discrepancy
therefore changes the output by at most
\begin{align*}
 \frac\rho2\sqrt{8MK}\,(2\sqrt K)^{t-2}
       2\sqrt{K/M}
 &=\rho\,2^{t-1/2}K^{t/2}\\
 &\leq\eta\,2^{-3/2}\leq\frac\eta2,
 \qquad 2\leq t\leq L-1.
\end{align*}
There are $L-1$ sampling stages, contributing at most $1/24$ in total.
Adding the truncation error proves the approximation bound in
\eqref{eq:compression-conclusion}; event $C$ proves its norm bound.

If the original network has strict margin $1/2$, the sampled network
has strict margin $5/12$. Multiplying its linear readout by $3/2$
restores strict margin $1/2$ and multiplies the weighted norm by at
most $9/4$. At width $M$ the resulting entrywise bounds are
\[
 |W^1_{aj}|\leq\sqrt{9KM/d},\qquad
 |W^s_{ab}|\leq3\sqrt K\ (1<s<L),\qquad
 |W^L_b|\leq3\sqrt{K/M}.
\]
Each is at most $M$ because $d\geq1$ and $M\geq9K\geq9$.
The definition of $C_L$ now applies. Finally, substituting
\eqref{eq:compression-parameters} gives
$\rho^{-4}=12^4(L-1)^4 16^L K^{2(L-1)}$, which establishes the stated
polynomial dependence at fixed depth.
\end{proof}


\begin{thebibliography}{9}
\bibitem{notes} Dmitry Vaintrob and Kaarel H\"anni.
\emph{A Generalization Result for Overparametrized Neural Net Bayesian Learning}.
Lecture notes, 2026.
\url{https://kaarelh.github.io/ai/nnbayes-generalization.pdf}.
\bibitem{fable} Claude Fable 5.
\emph{The circuit prior across the whole strong prior regime: cloning and subsampling at every $\gamma>1$}.
Draft dated August 27, 2026, prepared from an outline by Dmitry Vaintrob and Kaarel H\"anni; supplied by the authors.
\bibitem{critical} Dmitry Vaintrob, Kaarel H\"anni, and OpenAI Codex.
\emph{Bayesian circuit learning at mean-field scaling: a variational principle and width-independent generalization}.
Companion manuscript, October 2026.
\bibitem{subcritical} Dmitry Vaintrob, Kaarel H\"anni, and OpenAI Codex.
\emph{The fixed-kernel regime of a margin-conditioned Bayesian neural-network prior}.
Companion manuscript, October 2026.
\bibitem{borell} Christer Borell.
The Brunn--Minkowski inequality in Gauss space.
\emph{Inventiones Mathematicae}, 30:207--216, 1975.
\bibitem{ledoux} Michel Ledoux.
\emph{The Concentration of Measure Phenomenon}.
Mathematical Surveys and Monographs 89, American Mathematical Society, 2001.
\bibitem{davidson} Kenneth R. Davidson and Stanis\l aw J. Szarek.
Local operator theory, random matrices and Banach spaces.
In \emph{Handbook of the Geometry of Banach Spaces}, Vol. 1, pp. 317--366, Elsevier, 2001.
\bibitem{hoeffding} Wassily Hoeffding.
Probability inequalities for sums of bounded random variables.
\emph{Journal of the American Statistical Association}, 58(301):13--30, 1963.
\end{thebibliography}
\end{document}
